% !TEX root = ../paper.tex

\section{Additional Permit Evidence}
\label{sec:appendix_permits}

This appendix reports additional permit evidence that complements the descriptive discussion in Section~\ref{sec:data} and the event study in Section~\ref{sec:empirical_results_permits}.

\subsection{Permit Event-Study Estimates}

Table~\ref{tab:permit_event_study_did} reports pooled estimates for years 0 through 5 after redistricting.
The decline is concentrated in high-discretion applications.
The corresponding estimate for low-discretion applications is small and statistically indistinguishable from zero.

\begin{table}[htbp]
    \centering
    \caption{Alderman Stringency and Permit Applications After Redistricting}
    \label{tab:permit_event_study_did}
    \input{../tasks/run_event_study_permit/output/did_table_permit_2015_application_frozen2014_stable_preperiod_controls_500ft_clust_ward_pair.tex}

    \footnotesize
    \textit{Notes:} PPML estimates use census-block-year observations from 2010--2020 within 500ft of a ward boundary.
    The sample retains blocks whose 2014 origin and destination aldermen remained in office in 2015.
    Stringency scores are estimated using permits filed from 2006 through 2014.
    Each model includes block fixed effects, ward-pair $\times$ year fixed effects, and the outcome's pre-period volume and zero-volume indicator interacted with year.
    Fixed-effect PPML drops blocks with no outcome in any sample year and ward-pair-year groups with no outcome.
    Standard errors, in parentheses, are clustered by ward pair.
    $^{*}p<0.10$, $^{**}p<0.05$, $^{***}p<0.01$.
\end{table}

Figure~\ref{fig:permit_event_study_low_discretion} shows the full event-study path for low-discretion applications.
Unlike high-discretion applications, they do not fall after blocks move to more stringent aldermen.

\begin{figure}[htbp]
    \centering
    \includegraphics[width=0.7\textwidth]{../tasks/run_event_study_permit/output/event_study_permit_2015_low_discretion_nosigns_application_frozen2014_stable_preperiod_controls_500ft_clust_ward_pair.pdf}
    \caption{Placebo Event Study: Low-Discretion Permit Applications}
    \figurenotes{The specification and sample match Figure~\ref{fig:permit_event_study}, but the outcome is low-discretion permit applications excluding signs.
    Points report $\exp(\hat\beta_\tau)-1$; shaded areas are 95\% confidence intervals based on standard errors clustered by ward pair.
    The pooled post-period coefficient is $-0.017$ (SE $=0.048$; $p=0.731$).
    Joint pre-trend test $p=0.090$.}
    \label{fig:permit_event_study_low_discretion}
\end{figure}

\subsection{Permit Summary Statistics}
\label{app:permit_summary_stats}

Table~\ref{tab:permit_summary_stats} compares the high- and low-discretion permit groups used throughout the paper.
High-discretion permits are permits for projects where aldermen are most likely to matter: new construction, major renovation, demolition, porch construction, and reinstated permits.
They are much slower than lower-discretion permits: the mean processing time is 57.3 days versus 26.2 days, and the median is 30 days versus 6 days.
They also display more cross-alderman variation in processing times.

\begin{table}[htbp]
    \centering
    \caption{Summary Statistics for High- and Low-Discretion Permits}
    \label{tab:permit_summary_stats}
    \input{../tasks/permit_summary_stats/output/permit_processing_time_high_vs_low_alderman_summary.tex}

    \footnotesize
    \textit{Notes:} Sample uses permit applications through December 2022.
    High-discretion permits are new construction, renovation/alteration, wrecking/demolition, porch construction, and reinstated permits.
    Low-discretion permits are all remaining permit types excluding signs.
    The alderman-level interquartile range is computed from the cross-section of alderman-specific mean processing times within each permit group.
\end{table}

\subsection{Permit Volume and Processing Time}

Table~\ref{tab:permit_volume_correlation} reports the correlation between mean processing time and permit volume across ward-month cells.
The negative correlation for high-discretion permits is consistent with political and regulatory friction reducing throughput where discretionary review matters most, while the positive correlation for low-discretion permits suggests routine administrative scale rather than aldermanic discretion.

\begin{table}[htbp]
    \centering
    \caption{Processing Time and Permit Volume Correlation}
    \label{tab:permit_volume_correlation}
    \input{../tasks/permit_summary_stats/output/permit_processing_time_volume_correlation.tex}

    \footnotesize
    \textit{Notes:} Sample uses permit applications through December 2022.
    Entries report the Pearson correlation between mean processing time and permit volume across ward-month cells for each permit group.
\end{table}
